% در این فایل، دستورها و تنظیمات مورد نیاز، آورده شده است.
%-------------------------------------------------------------------------------------------------------------------

% در ورژن جدید زی‌پرشین برای تایپ متن‌های ریاضی، این سه بسته، حتماً باید فراخوانی شود
\usepackage{amsthm,amssymb,amsmath,mathrsfs,graphicx}
\usepackage{array}
\usepackage{booktabs,tabularx}
\usepackage[all]{xy}
% بسته‌ای برای تنطیم حاشیه‌های بالا، پایین، چپ و راست صفحه
\usepackage[top=40mm, bottom=25mm, left=25mm, right=35mm]{geometry}
% بسته‌‌ای برای ظاهر شدن شکل‌ها و تصاویر متن
\usepackage{graphicx,caption,subcaption}
\usepackage{adjustbox}
\usepackage{comment}
\usepackage{fancyhdr}
\usepackage{fontspec}
\usepackage{framed}
\usepackage{lastpage}
\usepackage{float}
\usepackage[table]{xcolor}
\usepackage{colortbl}
% بسته‌ و دستوراتی برای ایجاد لینک‌های رنگی با امکان جهش
\usepackage[pagebackref=false,colorlinks,linkcolor=blue,citecolor=magenta]{hyperref}
\hypersetup{urlcolor=blue}

\usepackage[nottoc]{tocbibind}
\usepackage{bidipoem}
% بسته‌ای برای رسم کادر
% بسته‌‌ای برای چاپ شدن خودکار تعداد صفحات در صفحه «معرفی پایان‌نامه»

% چنانچه قصد پرینت گرفتن نوشته خود را دارید، خط بالا را غیرفعال و  از دستور زیر استفاده کنید چون در صورت استفاده از دستور زیر‌‌،
% لینک‌ها به رنگ سیاه ظاهر خواهند شد که برای پرینت گرفتن، مناسب‌تر است
%\usepackage[pagebackref=false]{hyperref}
% بسته‌ لازم برای تنظیم سربرگ‌ها
% بسته‌ای برای ظاهر شدن «مراجع» و «نمایه» در فهرست مطالب

% دستورات مربوط به ایجاد نمایه
\usepackage{makeidx}
\makeindex
%%%%%%%%%%%%%%%%%%%%%%%%%%
% فراخوانی بسته زی‌پرشین و تعریف قلم فارسی و انگلیسی
\usepackage{xepersian}
\settextfont[Path=Fonts/,Scale=1]{XB Niloofar.ttf}

%%%%%%%%%%%%%%%%%%%%%%%%%%
% چنانچه می‌خواهید اعداد در فرمول‌ها، انگلیسی باشد، خط زیر را غیرفعال کنید
%\setdigitfont[Path=Fonts/,Scale=1]{Persian Modern.ttf}
%%%%%%%%%%%%%%%%%%%%%%%%%%
% تعریف قلم‌های فارسی و انگلیسی اضافی برای استفاده در بعضی از قسمت‌های متن
\defpersianfont\nastaliq[Path=Fonts/,Scale=2]{IranNastaliq.ttf}
\defpersianfont\chapternumber[Path=Fonts/,Scale=3]{XB Niloofar.ttf}
\defpersianfont\naz[Path=Fonts/,Scale=0.9]{B Nazanin.TTF}
\defpersianfont\nazz[Path=Fonts/,Scale=1]{B Nazanin.TTF}
\defpersianfont\nazzz[Path=Fonts/,Scale=1.1]{B Nazanin.TTF}
\defpersianfont\nazzzz[Path=Fonts/,Scale=1.4]{B Nazanin.TTF}
\defpersianfont\nazzzzz[Path=Fonts/,Scale=1.7]{B Nazanin.TTF}
\defpersianfont\Titr[Path=Fonts/,Scale=1.4]{Titr.ttf}
\deflatinfont\Tim[Path=Fonts/,Scale=.9]{Times New Roman.ttf}
\deflatinfont\Timm[Path=Fonts/,Scale=1]{Times New Roman.ttf}
\deflatinfont\Timmm[Path=Fonts/,Scale=1.1]{Times New Roman.ttf}
\deflatinfont\Timmmm[Path=Fonts/,Scale=1.2]{Times New Roman.ttf}
\deflatinfont\Timmmmm[Path=Fonts/,Scale=1.3]{Times New Roman.ttf}
\deflatinfont\Timmmmmm[Path=Fonts/,Scale=1.4]{Times New Roman.ttf}

%%%%%%%%%%%%%%%%%%%%%%%%%%
% دستوری برای حذف کلمه «چکیده»
%\renewcommand{\abstractname}{}
% دستوری برای حذف کلمه «abstract»
\renewcommand{\latinabstract}{}
% دستوری برای تغییر نام کلمه «اثبات» به «برهان»
\renewcommand\proofname{\textbf{برهان}}
% دستوری برای تغییر نام کلمه «کتاب‌نامه» به «مراجع»
\renewcommand{\bibname}{مراجع}
\renewcommand{\indexname}{فهرست نمادها}
% دستوری برای تعریف واژه‌نامه انگلیسی به فارسی
\newcommand\persiangloss[2]{#1\dotfill\lr{#2}\\}
% دستوری برای تعریف واژه‌نامه فارسی به انگلیسی
\newcommand\englishgloss[2]{#2\dotfill\lr{#1}\\}
% تعریف دستور جدید «\پ» برای خلاصه‌نویسی جهت نوشتن عبارت «پروژه/پایان‌نامه/رساله»
\newcommand{\پ}{پروژه/پایان‌نامه/رساله }
%%%%%%%%%%%%%%%%%%%%%%%%%%
% تعریف و نحوه ظاهر شدن عنوان قضیه‌ها، تعریف‌ها، مثال‌ها و ...
\theoremstyle{definition}
\newtheorem{definition}{تعریف}[chapter]
\theoremstyle{plain}
%\newtheorem{s}[definition]{}
\newtheorem{theorem}[definition]{قضیه}
\newtheorem{lemma}[definition]{لم}
\newtheorem{proposition}[definition]{گزاره}
\newtheorem{corollary}[definition]{نتیجه}
\newtheorem{remark}[definition]{تبصره}
\theoremstyle{definition}
\newtheorem{example}[definition]{مثال}
\newtheorem{notation}[definition]{نمادگذاری}
\newtheorem{definitiont}[definition]{تعریف و نمادگذاری}
\newtheorem{obs}[definition]{ملاحظه}
%%%%%%%%%%%%%%%%%%%%%%%%%%%%

% دستورهایی برای سفارشی کردن سربرگ صفحات
\csname@twosidetrue\endcsname
\pagestyle{fancy}
\setlength{\headheight}{15pt}
\emergencystretch=2em
\fancyhf{}
\fancyhead[OR,ER]{\thepage}
\fancyhead[OL]{\small\rightmark}
\fancyhead[EL]{\small\leftmark}
\renewcommand{\chaptermark}[1]{%
\markboth{\thechapter.\ #1}{}}
%%%%%%%%%%%%%%%%%%%%%%%%%%%%%

\newcommand{\rt}{\rightarrow}
\newcommand{\lrt}{\longrightarrow}
\newcommand{\lto}{\longrightarrow}
\newcommand{\st}{\stackrel}

\newcommand{\X}{\textbf{X}}
\newcommand{\Y}{\textbf{Y}}
\newcommand{\Z}{\mathbb{Z}}
\newcommand{\W}{\textbf{W}}
\newcommand{\PP}{\textbf{P}}
\newcommand{\II}{\textbf{I}}
\newcommand{\QQ}{\textbf{Q}}
\newcommand{\bb}{{\rm b}}
\newcommand{\CP}{\mathcal{P}}
\newcommand{\F}{\mathbf{F}}
\newcommand{\BF}{\mathbf{F}}
\newcommand{\CQ}{\mathcal{Q}}
\newcommand{\CD}{\mathcal{D}}
\newcommand{\CCM}{\mathcal{M}}
\newcommand{\Lf}{\mathcal{L}f}
\newcommand{\Sf}{\mathcal{S}f}

\newcommand{\Ga}{\Gamma}
\newcommand{\La}{\Lambda}
\newcommand{\om}{\omega}
\newcommand{\al}{\alpha}
\newcommand{\la}{\lambda}
\newcommand{\Tr}{{\rm{Tr}}}
\newcommand{\Dtr}{{\mathbf{D}\rm{Tr}}}
\newcommand{\DD}{{\rm{D}}}
\newcommand{\J}{\rm{J}}
\newcommand{\tc}{\bigotimes_{\mathcal{C}}}

\newcommand{\cm}{\mathbf{C}(\rm{mod} \Lambda)}
\newcommand{\cmo}{\mathbf{C}(\rm{mod} \Lambda^{\rm{op}})}
\newcommand{\cmb}{\mathbf{C}^{\rm{b}}(\rm{mod} \Lambda)}
\newcommand{\mmod}{{\rm mod}\mbox{-}}
\newcommand{\Mod}{{\rm Mod}\mbox{-}}
\newcommand{\prj}{{\rm prj}\mbox{-}}
\newcommand{\Gprj}{{\rm Gprj}\mbox{-}}
\newcommand{\Ginj}{{\rm Ginj}\mbox{-}}
\newcommand{\rad}{{\rm rad}~}
\newcommand{\inj}{{\rm inj}\mbox{-}}
\newcommand{\soc}{{\rm soc}}
\newcommand{\Tor}{{\rm Tor}}
\newcommand{\Rep}{{\rm Rep}}
\newcommand{\rep}{{\rm rep}}
\newcommand{\Inj}{{\rm Inj}\mbox{-}}
\newcommand{\GInj}{{\rm GInj}\mbox{-}}
\newcommand{\Prj}{{\rm Prj}\mbox{-}}
\newcommand{\GPrj}{{\rm GPrj}\mbox{-}}
\newcommand{\Sum}{{\rm Sum}\mbox{-}}
\newcommand{\summ}{{\rm sum}\mbox{-}}
\newcommand{\Flat}{{\rm Flat}\mbox{-}}
\newcommand{\Vecc}{{\rm Vec}}
\newcommand{\vecc}{{\rm vec}}
\newcommand{\gr}{{\rm gr}\mbox{-}}
\newcommand{\add}{{\rm add}\mbox{-}}


\newcommand{\End}{{\rm{End}}}
\newcommand{\lan}{\langle}
\newcommand{\ran}{\rangle}

\newcommand{\BC}{\mathbb{C} }
\newcommand{\R}{\mathbb{R} }
\newcommand{\BD}{\mathbf{D}}
\newcommand{\BZ}{\mathbf{Z}}
\newcommand{\D}{\mathbb{D} }
\newcommand{\K}{\mathbb{K} }
\newcommand{\CM}{{\rm CM}\mbox{-}}
\newcommand{\CF}{\mathcal{F}}
\newcommand{\CI}{\mathcal{I}}
\newcommand{\CT}{\mathcal{T}}
\newcommand{\CV}{\mathcal{V}}
\newcommand{\CS}{\mathcal{S}}
\newcommand{\CY}{\mathcal{Y}}
\newcommand{\CX}{\mathcal{X}}
\newcommand{\BI}{\mathbf{I}}


\newcommand{\N}{\mathbb{N} }
\newcommand{\Q}{\mathbf{Q} }
%\newcommand{\Z}{\mathbb{Z} }

\newcommand{\CA}{\mathcal{A} }
\newcommand{\CC}{\mathcal{C} }
\newcommand{\BX}{\mathbf{X}}
\newcommand{\BY}{\mathbf{Y}}

\newcommand{\im}{{\rm{Im}}}
\newcommand{\op}{{\rm{op}}}
\newcommand{\inc}{{\rm{inc}}}
\newcommand{\can}{{\rm{can}}}
\newcommand{\perf}{{\rm perf}}
\newcommand{\coperf}{{\rm coperf}}

\newcommand{\Coker}{{\rm{Coker}}}
\newcommand{\Ker}{{\rm{Ker}}}

\newcommand{\Hom}{{\rm{Hom}}}
\newcommand{\Ext}{{\rm{Ext}}}

%% boldface letters
\def\bfm#1{\protect{\makebox{\boldmath $#1$}}}
\def\b {\bfm{b}}
\def\a {\bfm{a}}
\def\aalpha {\bfm{\alpha}}
\def\bbeta {\bfm{\beta}}
\def\eell {\bfm{\ell}}
\def\llambda {\bfm{\lambda}}
\def\pphi {\bfm{\phi}}
\def\f {\bfm{f}}
\def\p {\bfm{p}}
\def\q {\bfm{q}}


%%%%%%%%%%%%%%%%%%%%%%
% دستورهایی برای سفارشی کردن صفحات اول فصل‌ها
\makeatletter
\newcommand\mycustomraggedright{%
 \if@RTL\raggedleft%
 \else\raggedright%
 \fi}
\def\@part[#1]#2{%
\ifnum \c@secnumdepth >-2\relax
\refstepcounter{part}%
\addcontentsline{toc}{part}{\thepart\hspace{1em}#1}%
\else
\addcontentsline{toc}{part}{#1}%
\fi
\markboth{}{}%
{\centering
\interlinepenalty \@M
\ifnum \c@secnumdepth >-2\relax
 \huge\bfseries \partname\nobreakspace\thepart
\par
\vskip 20\p@
\fi
\Huge\bfseries #2\par}%
\@endpart}
\def\@makechapterhead#1{%
\vspace*{-30\p@}%
{\parindent \z@ \mycustomraggedright %\@mycustomfont
\ifnum \c@secnumdepth >\m@ne
\if@mainmatter
\vskip 100\p@
\huge\bfseries \@chapapp\space {{\chapternumber}\thechapter}
\par\nobreak
\vskip 40\p@
\fi
\fi
\interlinepenalty\@M
\Huge \bfseries #1\par\nobreak
\vskip 70\p@
}}
\renewcommand\tableofcontents{%
    \if@twocolumn
      \@restonecoltrue\onecolumn
    \else
      \@restonecolfalse
    \fi
    \chapter*{\vspace*{2cm}\contentsname
        \@mkboth{\MakeUppercase\contentsname}{\MakeUppercase\contentsname}}%
    \@starttoc{toc}%
    \if@restonecol\twocolumn\fi
}
\makeatother 
